\begin{abstract}
We implement and train an RGB-conditioned Diffusion Policy for ManiSkill's
\texttt{PushT-v1} task. The reportable experiment uses 100 successful, controller-native
demonstrations replayed with temporally aligned RGB observations, two observation frames, a
16-step prediction horizon, one-step receding-horizon execution, and 50,000 optimization steps.
The selected EMA checkpoint is fixed at iteration 15,000, where a 20-episode diagnostic first
reaches 30\% \successonce. On the required independent evaluation of 100 rollouts at each of
three environment seeds, it obtains \textbf{\timean\% $\pm$ \tistd{} percentage points}
(\tipooled{} pooled successes).

The 1,000-step moving-average diffusion loss decreases from 0.0533 at iteration 11,000 to 0.0183
at 50,000. The preserved Runpod continuation from 10,000 to 50,000 takes 2 h 19 min 2 s across
two restart-safe segments; the uninterrupted 15,000--50,000 segment takes 2 h 2 min 35 s on one
NVIDIA A40 and peaks at 7,603 MiB (7.42 GiB) sampled VRAM. We also document the principal
debugging result: an earlier GPU replay paired actions with a one-step-shifted environment state.
Correcting the replay to restore $s_{t+1}$ after action $a_t$, matching the original 1,024-way
simulation parallelism, and retraining from a valid checkpoint transformed a zero-success
pipeline into a measurable Push-T baseline. The final result remains limited by one training
seed and a noisy checkpoint-selection diagnostic; these limitations motivate the targeted
failure analysis in T-II. T-II tooling for pose-conditioned resets, trajectory capture, and
failure screening was implemented and tested, but the required discovery and targeted rollout
experiments did not complete because replacement GPU instances either lacked a working Vulkan
graphics path or were unavailable. Consequently, this report does not claim two validated failure
modes, targeted success rates, or representative T-II videos. T-III and T-IV were also not
completed; we instead state the intended LLM reward-generation and bounded residual-PPO methods,
including the validation needed to avoid reward hacking and distributional regressions.
\end{abstract}
